\chapter{Training prerequisite map: 무엇이 실제로 바뀌는가}
\label{STC-S03}

Sleep-time compute를 평가하려면 ``learning''을 하나의 동사로 쓰지 않고 data, objective, trainable
state, optimizer, validation, serving artifact로 분해해야 한다. 이 장은 별도 Training Background와
Study 사이의 navigation layer다.

\section{첫 용어 지도}

\begin{table}[htbp]
\centering
\caption{Study에서 반복 사용할 training 개념과 Background 링크}
\label{tab:study-training-map}
\small
\begin{tabularx}{\textwidth}{p{0.23\textwidth} X X}
\toprule
개념 & 최소 의미 & STC에서 묻는 질문 \\
\midrule
Fine-tuning\bgref{fine-tuning} / continued pre-training\bgref{continued-pretraining} & 좁은 목적 또는 추가 corpus로 배포 모델을 갱신 & base 전체인가, delta인가? \\
LoRA\bgref{lora} / adapter\bgref{adapter} & 작은 trainable state로 behavior를 조정 & user/domain별 artifact 수와 merge는? \\
Distillation\bgref{distillation} & teacher behavior를 student로 압축 & tail·uncertainty·provenance가 보존되는가? \\
Replay buffer\bgref{replay-buffer} / augmentation\bgref{data-augmentation} & 과거·변형 data를 new data와 섞음 & retention과 recursive-data depth는? \\
Continual learning\bgref{continual-learning} & stream을 배우며 이전 능력을 유지 & multi-cycle forgetting과 capacity knee는? \\
RLHF/RL\bgref{rlhf} & delayed/preference signal로 policy를 조정 & model이 아니라 trigger/router를 학습할 것인가? \\
Meta-learning\bgref{meta-learning} & 빠른 adaptation이 가능하도록 학습 & user sleep inner loop가 안정적인가? \\
Test-time training\bgref{test-time-training} & test stream에서 state를 즉시 갱신 & current response path인가 deferred path인가? \\
Model editing\bgref{model-editing} / unlearning\bgref{unlearning} & 좁은 수정 또는 영향 제거 & locality, lineage, rollback 단위는? \\
\bottomrule
\end{tabularx}
\end{table}

\section{Operator record}

각 연구 방법을 다음 record로 정규화한다.

\begin{equation}
R_{op}=\{D, T, \theta_{train}, \mathcal{L}/r, U, C, V, A\},
\end{equation}

$D$는 wake-derived data와 replay, $T$는 teacher/source, $\theta_{train}$은 바뀌는 state,
$\mathcal{L}/r$은 loss 또는 reward, $U$는 update rule/optimizer, $C$는 capacity destination,
$V$는 validation slices, $A$는 published artifact와 rollback unit이다. 이 record가 비어 있으면
``모델이 기억을 공고화했다''는 설명을 system으로 환원할 수 없다.

\begin{table}[htbp]
\centering
\caption{대표 sleep operator를 training record로 펼친 예}
\label{tab:study-operator-record}
\scriptsize
\begin{tabularx}{\textwidth}{p{0.12\textwidth} p{0.13\textwidth} p{0.13\textwidth} p{0.14\textwidth} X X}
\toprule
Operator & Data/teacher & Trainable state & Objective/update & Capacity & Artifact/gate \\
\midrule
summary synthesis & episodes, LLM teacher & 없음 또는 prompt/model & fidelity, utility, compression & external text & versioned record, QA \\
temporal graph & event+time+entity & graph nodes/edges & merge, supersede, conflict & external graph & graph epoch, traversal test \\
LoRA sleep & new+replay & low-rank $A,B$ & task+retention loss, Adam & adapter bank & candidate delta, shadow \\
distillation & teacher outputs+real anchor & student/base delta & CE+KL+retention & model artifact & full release suite \\
RL controller & trajectories & controller $\phi$ & delayed utility-cost reward & policy model & constrained canary \\
model edit & fact+neighbors & localized weight delta & efficacy+locality & parametric patch & edit version, rollback \\
\bottomrule
\end{tabularx}
\end{table}

\section{Training cost는 FLOPs만이 아니다}

Full fine-tuning은 weights 외에 gradient, optimizer state, activation을 필요로 한다. LoRA는 trainable
parameter와 optimizer bytes를 줄이지만 base read, forward/backward activation, validation, version
management는 남는다. External transformation은 GPU training이 없을 수 있지만 teacher inference,
embedding, graph write, dedup, index rebuild와 evaluation을 낸다.

따라서 end-to-end sleep cost를

\begin{equation}
C_{sleep}=C_{ingest}+C_{transform}+C_{train}+C_{validate}+C_{publish}+C_{retain}
\end{equation}

로 분리한다. $C_{retain}$은 artifact가 존재하는 동안의 storage, migration, rebase, deletion debt다.
Kernel FLOPs만 비교하면 external path의 retrieval cost와 parametric path의 lifecycle cost를 서로 다른
회계 기준으로 비교하게 된다.

\section{Wake, sleep, long-term memory의 세 가지 계약}

Wake path는 response SLA와 isolation을, sleep path는 throughput·deadline·candidate safety를, long-term
memory plane은 persistence·scope·version·delete를 계약한다. 한 cluster에서 모두 실행할 수도 있지만
logical contract는 분리한다. Shared state는 explicit snapshot epoch를 가져야 wake가 training 중간값을
읽지 않는다.

\begin{keypoint}[논문을 읽는 최소 checklist]
Data는 어디서 왔는가? 어느 state가 얼마나 오래 바뀌는가? Old/new/safety를 무엇으로 검증했는가?
Capacity가 찰 때 무슨 state를 제거하는가? 후보와 serving state를 어떻게 전환하는가? 이 다섯 답이
없으면 accuracy gain을 production sleep system으로 해석하지 않는다.
\end{keypoint}
